\cotefile{123}
\scanpage{15}
\uline{Résumé} th. \unsure{de|du} conjugaison Kostant $+$ lettre à Kostant
% archive sheet number: 14
\scanpage{24}
\unsure{Diagramme|Diagrammes} d'espaces homogènes \unsure{reliés|liés} aux éléments de Kostant, de Coxeter et aux pts cycliques.
% archive sheet number: 23
\scanpage{26}
\margin{!!}
% slanted
\margin{\unsure{Valid.|Valab.|Vrais} \unsure{seulem$^{\mathrm{t}}$|seult|seulement} en Car. 0 (ou car. $p>h$ ?)} 
\diagram{\resizebox{\linewidth}{!}{$\begin{array}{@{}l@{\quad}l@{\quad}l@{}}
\begin{array}{@{}l@{}}\text{éléments de Kostant}\quad\unsure{\kappa|k}:\mu_h\longrightarrow G\\\text{éléments de Coxeter}\quad\unsure{r|\gamma}\in G(k)\end{array} & \left.\begin{array}{@{}l@{}}\mathrm{E\,Kost}\\\mathrm{E\,Cox}\end{array}\right\}\begin{array}{@{}l@{}}\text{kif-kif si car. 0}\\\text{et on se donne}\\\omega\in\mu_h^*(k)\end{array} & \begin{array}{@{}l@{}}\deleted{\text{\unsure{splittée}}}\\\deleted{\text{ext. de }Z\text{ par }T}\quad\unsure{\mathcal Z|Z|\mathfrak z}\mathbin{.}T\quad(\tfrac12\ \text{direct})\end{array} \\[1.2ex]
\begin{array}{@{}l@{}}\text{\unsure{Ssgroupes|Ss groupes|Ss-groupes} \inserted{de} Kostant}\quad\mu\hookrightarrow G\\\text{\unsure{Ss groupe|Ss-groupe|ss-groupes} de Coxeter}\quad H\hookrightarrow G\end{array} & \left.\begin{array}{@{}l@{}}\mathrm{SG\,Kost}\\\mathrm{SG\,Cox}\end{array}\right\}\begin{array}{@{}l@{}}\text{kif kif en}\\\text{car 0}\end{array} & \text{ext. de }(\deleted{\unsure{\mu}}\mathbb Z/h\mathbb Z)^*\ \text{par }\mathcal Z\mathbin{.}\deleted{\text{\unsure{par}}}T \\[1.2ex]
\left\{\begin{array}{@{}l@{}}\text{Couples de Kostant}\quad(\kappa:\mu_h\overset{\deleted{\unsure{T}}}{\hookrightarrow}G,T')\\\quad\inserted{\text{ou au choix }(\kappa:\mu\hookrightarrow G,\underset{\substack{\uparrow\\\unsure{\mathfrak g(1)^{\mathrm{reg}}|g(1)^{\mathrm{reg}}|\hat{\mathfrak g}(1)^{\mathrm{reg}}}}}{B})}\\\text{Couple de Coxeter}\quad(\gamma,T)\end{array}\right. & \left.\begin{array}{@{}l@{}}\mathrm{C\,Kost}\\\mathrm{C\,Cox}\end{array}\right\}\begin{array}{@{}l@{}}\text{kif kif si car. 0}\\\text{\unsure{si on se donne}}\\\omega\in\mu_h^*(k)\end{array} & \mu_h\times\mathcal Z \\[1.2ex]
\begin{array}{@{}l@{}}\text{id}\ \left\{\begin{array}{@{}l@{}}\text{précouple de \(\prime\prime\) Kostant}\quad(\mu\subset G,T')\\\text{précouple de \(\prime\prime\) Coxeter}\quad(\unsure{H|\mu}\subset G,T)\end{array}\right.\\\qquad\text{couples de tores en opposition}\end{array} & \begin{array}{@{}l@{}}\deleted{\mathrm{PC\,Kost}\quad\text{kif kif en car. 0}}\\\deleted{\mathrm{PC\,Cox}}\\{\unsure{\mathrm{Torapp}|\mathrm{Toropp}|\mathrm{Toapp}}}^{*}\end{array} & \begin{array}{@{}l@{}}\text{ext. de }(\mathbb Z/h\mathbb Z)^*\ \text{par }\mu_h\times\mathcal Z,\\\text{(à préciser)}\end{array} \\[1.2ex]
\begin{array}{@{}l@{}}\text{droites cycliques}\quad D\\\text{\deleted{\unsure{Couples|Classe} \illegible{1}}\inserted{Classe} de Coxeter}\quad(T,\underset{\substack{\rotatebox{90}{$\in$}\\W(k)}}{\dot\gamma})\\\text{élément \deleted{cyclique}\inserted{\deleted{\unsure{normé}} cyclique normalisé} (dépend du }{}_{\circ}\! /\end{array} & \begin{array}{@{}l@{}}\left.\begin{array}{@{}l@{}}\mathrm{DCy}\\\unsure{\mathrm{Cl\,Cox}|\mathrm{CE\,Cox}}\end{array}\right\}\begin{array}{@{}l@{}}\text{kif kif si car. 0}\\\text{et on donne}\\\omega\in\mu_h^*(k)\end{array}\\\mathrm{ECy}_{a}\end{array} & \begin{array}{@{}l@{}}\mu_h\mathbin{.}T'\quad(\tfrac12\ \text{direct})\\T'\end{array} \\[1.2ex]
\begin{array}{@{}l@{}}\text{\deleted{triple}\inserted{couple} de Kostant\inserted{-Killing}}\quad((B,T),T')\ \unsure{\text{ou}|\sim}\ ((B,T),D)\\\qquad(\text{tel que }T'\text{ normalisé}\\\qquad\text{par }\mu_h\xrightarrow{\kappa}T\text{ défini par }B\supset T)\end{array} & \deleted{\unsure{C|E}}\ \mathrm{Kost\,Kill} & \mu_h \\[1.2ex]
\begin{array}{@{}l@{}}\text{\deleted{triple}\inserted{couple} de \deleted{Coxeter}\inserted{Kostant normalisé}}\quad(T,\deleted{\unsure{D}}x)\\\qquad\text{\deleted{droite \unsure{critique|cyclique} définissant}\inserted{point cyclique normalisé}}\\\qquad\text{un tore }T'\text{ en opp. avec }T,\\\qquad\text{i.e.}\end{array} & \deleted{\unsure{\mathrm{Pr\,Cox}|\mathrm{PC\,Cox}|\mathrm{PCox}}}\ \mathrm{C\,Kost\,norm}_{a} & \mathcal Z \\[1.2ex]
\text{Couple de Kostant-Killing normalisé}\quad((T,B),x) & \mathrm{Kost\,Kill\,norm}_{a} & 1
\end{array}$}}

\noindent $\mathord{}_{\circ}\! /$ choix d'un $a\in\Delta^*(k)$, où $\Delta$ est le quotient commun aux droites cycliques par $\mu_h$ (\uline{NB} modulo signe, il y a un choix privilégié \deleted{\unsure{dans}} d'un $x$ \deleted{$\unsure{g|y|\mathfrak z}$}, \unsure{des|deux} \unsure{versions|racines} \unsure{sur|par} ext. quadratique)
% archive sheet number: 25
\scanpage{27}
\diagram{\resizebox{\linewidth}{!}{\begin{tikzcd}[ampersand replacement=\&,row sep=1.4em,column sep=0.8em,row sep=.3em,column sep=.75em]
 \&  \&  \& 0\arrow[dash,to=2-3]\arrow[dash,to=2-5] \&  \&  \\
 \&  \& \unsure{\mathcal Z|Z|\mathfrak z}\arrow[dash,to=3-4] \&  \& \mu\arrow[dash,to=3-4] \&  \\
 \&  \&  \& \mu\times\mathcal Z\arrow[dash,to=4-3,"{T/\mu}"]\arrow[dash,to=4-5]\arrow[dash,to=5-4,"{(\mathbb Z/h\mathbb Z)^*}"] \&  \&  \\
 \&  \& \mathcal Z\mathbin{.}T\arrow[dash,to=6-2,"{(\mathbb Z/h\mathbb Z)^*}"] \&  \& \mu\mathbin{.}T'\arrow[dash,to=6-6] \&  \\
 \&  \&  \& N\cap N'\arrow[dash,to=6-2]\arrow[dash,to=6-6] \&  \&  \\
 \& N(\unsure{{}_{t}\mu|{}^{t}\mu|\mu})\arrow[dash,to=7-1] \&  \&  \&  \& N' \\
N \&  \&  \&  \&  \& 
\end{tikzcd}}}

\diagram{\resizebox{\linewidth}{!}{\begin{tikzcd}[ampersand replacement=\&,row sep=1.4em,column sep=0.8em,row sep=1em,column sep=.5em,execute at end picture={\draw[thin,rounded corners=5pt] ([xshift=-1em,yshift=1.3em]\tikzcdmatrixname-1-4.north west) -- ([xshift=-.8em,yshift=.7em]\tikzcdmatrixname-2-3.north west) -- ([xshift=-.9em,yshift=-.8em]\tikzcdmatrixname-2-3.south west) -- ([xshift=.6em]\tikzcdmatrixname-4-3.east) -- ([xshift=-.8em,yshift=-.9em]\tikzcdmatrixname-7-4.south west) -- ([xshift=.5em,yshift=-1.1em]\tikzcdmatrixname-7-4.south -| \tikzcdmatrixname-6-5.east) -- ([xshift=1em]\tikzcdmatrixname-4-6.east) .. controls +(2em,2em) and +(2em,0em) .. ([xshift=1em,yshift=1.3em]\tikzcdmatrixname-1-4.north east);\draw[thin] ([yshift=-1em]\tikzcdmatrixname-7-4.south) -- (\tikzcdmatrixname-9-4.north);}]
 \&  \&  \& \begin{gathered}\mathrm{Kost{-}Kill.norm}_{a}\\{\scriptstyle \deleted{\text{\unsure{paires} }((T,B),(T',B')),\ \text{\unsure{avec} }T,T'\text{ en opposition,}}}\\{\scriptstyle \deleted{\text{\unsure{une classe} \illegible{3} \unsure{de} }\unsure{W'/\mu}}}\end{gathered}\arrow[dash,to=2-3]\arrow[dash,to=2-5] \&  \&  \&  \\
 \&  \& \begin{gathered}\mathrm{C\,Kost\,norm}_{a}\\{\scriptstyle \deleted{\text{paires }(T,(T',B')),\text{ avec}}}\\{\scriptstyle \deleted{T\text{ en opposition avec }T',}}\\{\scriptstyle \deleted{B'\text{ dans une classe donnée}}}\\{\scriptstyle \deleted{\text{de }W'\!/\!(\text{\unsure{ss-gpe|ss-gr.} extension}}}\\{\scriptstyle \deleted{\text{de }(\mathbb Z/h\mathbb Z)^*\text{ par }\mu)}}\end{gathered}\arrow[dash,to=3-4] \&  \& \begin{gathered}\mathrm{Kost\,Kill}\\{\scriptstyle \text{paire }((T,B),T'),\text{ telle que}}\\{\scriptstyle (\unsure{\kappa|k},T')\text{ soit une paire}}\\{\scriptstyle \text{de Kost., }\kappa\text{ étant l'élément}}\\{\scriptstyle \text{de Kostant défini par le}}\\{\scriptstyle \text{couple de Killing}}\end{gathered}\arrow[dash,to=3-4] \&  \&  \\
 \&  \&  \& \mathrm{C\,Kos}\arrow[dash,to=4-3]\arrow[dash,bend right=11,to=7-4]\arrow[dash,to=5-4,"{\wr^{*}}"]\arrow[dash,to=4-6] \&  \&  \&  \\
 \&  \& \mathrm{E\,Kos}\arrow[dash,to=5-2]\arrow[dash,to=6-3,"{\wr^{**}}"] \&  \&  \& \mathrm{DCy}\arrow[dash,to=6-5,"{\wr^{**}}"]\arrow[dash,to=6-7] \&  \\
 \& \mathrm{SG\,Kos}\arrow[dash,to=7-1]\arrow[dash,to=7-2,"{\wr^{*}}"] \&  \& \mathrm{C\,Cox}\arrow[dash,to=6-3,"{\unsure{1|\vert}}"]\arrow[dash,to=7-4,"{*}"]\arrow[dash,to=6-5] \&  \&  \&  \\
 \&  \& \mathrm{E\,Cox}\arrow[dash,to=7-2] \&  \& \mathrm{PC\,Cox}\arrow[dash,to=7-7] \&  \& \mathrm{Tor}\arrow[phantom,sloped,to=7-7,"{\Vert}"] \\
\mathrm{Tor}\arrow[phantom,sloped,to=8-2,"{\Vert}"] \& \mathrm{SG\,Cox}\arrow[dash,to=8-2]\arrow[dash,to=7-4,"{*}"] \&  \& {\unsure{\mathrm{Torapp}|\mathrm{Toropp}|\mathrm{Toapp}}}^*\arrow[dash,to=7-7,"{*}"] \&  \&  \& \mathrm{Tor} \\
 \& \mathrm{Tor} \&  \&  \&  \&  \&  \\
 \&  \&  \& \substack{\text{leur existence dans}\\\text{cas quasi-déployé}\\\text{n'est pas claire.}} \&  \&  \& 
\end{tikzcd}}}

\margin{$**$ Isom défini si on travaille sur $\mu_h^*$\par $*$ $\mathrm{Isom}^{\circ}$ défini si $(h,p)=1$ ou \unsure{entier|entière$^{\mathrm{t}}$|notion} défini dans ce cas}
% archive sheet number: 26
\scanpage{29}
couple de tores \inserted{$(T,T')$} en opposition

$\simeq$ Ss-gpe de Coxeter $\mu$, \deleted{\unsure{un ss-gpe}} de Cartan normalisé
\inserted{tel que $\mu$ opère \unsure{fidèl$^{\mathrm{t}}$|transit$^{\mathrm{t}}$|lib$^{\mathrm{t}}$|fidèlement} sur une droite \unsure{convenable|can.|cyclique} de Cartan (\unsure{lire|lie|Lie} \unsure{gpe|alg|$\mathfrak g$})}

\unsure{or|et} le choix d'un isom $\mu\simeq\mu_h$ équivaut au choix d'un $D$ \unsure{\ldots|---}
\[
\left\{\begin{aligned}
&\text{\unsure{sous-gpe|sous-groupes} de Kostant-Coxeter}\\
&\text{Elément de Kostant}
\end{aligned}\right.
\]
Sous-g\unsure{p|pe}

\diagram{\resizebox{\linewidth}{!}{\begin{tikzcd}[ampersand replacement=\&,row sep=1.4em,column sep=0.8em,row sep=2.8em,column sep=2.3em]
T\arrow[phantom,to=2-2,"{\mathrm{coc.}}"]\arrow[dash,to=1-2,"{\mathcal Z}"]\arrow[dash,to=2-1,"{T/\mu}"] \& \unsure{\mathcal Z|Z|\mathfrak z}(\mu)\arrow[phantom,to=2-3,"{\mathrm{coc.}}"]\arrow[dash,to=1-3,"{(\mathbb Z/h\mathbb Z)^*}"]\arrow[dash,to=2-2,"{T/\mu}"] \& N(\mu)\arrow[dash,to=1-4]\arrow[dash,to=2-3,"{T/\mu}"] \& N\arrow[dash,to=2-4] \&  \\
T\cap N'\arrow[phantom,to=3-2,"{\mathrm{coc.}}"]\arrow[dash,to=2-2,"{\mathcal Z}"]\arrow[equals,to=3-1,"{0}"] \& \mathcal Z(\mu)\cap N'\arrow[dash,to=2-3,"{\deleted{\unsure{\mu}}(\mathbb Z/h\mathbb Z)^*}"]\arrow[equals,to=3-2,"{0}"] \& N(\mu)\cap N'\arrow[phantom,to=3-4,"{\mathrm{coc.}}"]\arrow[equals,to=2-4]\arrow[equals,to=3-3,"{(\mathbb Z/h\mathbb Z)^*}"] \& N\cap N'\arrow[dash,to=2-5,"{\unsure{-|\to}}"]\arrow[dash,to=3-4,"{(\mathbb Z/h\mathbb Z)^*}"] \& N'\arrow[dash,to=3-5] \\
\deleted{\unsure{D}}\ T\cap N(D)\arrow[phantom,to=4-2,"{\mathrm{coc.}}"]\arrow[dash,to=3-2,"{\mathcal Z}"]\arrow[dash,to=4-1,"{\mu}"] \& \mathcal Z(\mu)\cap N(D)\arrow[phantom,to=4-3,"{\mathrm{coc.}}"]\arrow[equals,to=3-3,"{0}"]\arrow[dash,to=4-2,"{\mu}"] \& N(\mu)\cap N(D)\arrow[phantom,to=4-4,"{\mathrm{coc.}}"]\arrow[equals,to=3-4]\arrow[dash,to=4-3,"{\mu}"] \& N\cap N(D)\,{}^{=\mu\times\mathcal Z}\arrow[phantom,to=4-5,"{\mathrm{coc.}}"]\arrow[dash,to=3-5,"{T'/\mathcal Z}"]\arrow[dash,to=4-4,"{\mu}"] \& N(D)=\mu\mathbin{.}T'\arrow[dash,to=4-5,"{\mu}"] \\
T\cap T'\arrow[dash,to=4-2,"{\mathcal Z}"]\arrow[equals,to=5-1] \& \mathcal Z(\mu)\cap T'\arrow[equals,to=4-3] \& N(\mu)\cap T'\arrow[equals,to=4-4] \& N\cap T'\arrow[dash,to=4-5,"{T'/\mathcal Z}"] \& T' \\
0 \&  \&  \&  \& \deleted{\unsure{T'}}\mathcal Z
\end{tikzcd}}}
\[
\deleted{\unsure{N(\mu)\cap N'}}=
\left\{\begin{aligned}
N\cap N(D)&=\mathcal Z(\mu)\cap N(D)\\
N\cap T'&=\mathcal Z(\mu)\cap T'\\
N\cap N'&=N(\mu)\cap N'\unsure{\mathclose{\}}|\,;}
\end{aligned}\right.
\]
\diagram{\resizebox{\linewidth}{!}{\begin{tikzcd}[ampersand replacement=\&,row sep=1.4em,column sep=0.8em,row sep=.6em]
T\arrow[dash,to=1-2]\arrow[dash,to=2-1] \& T\mathbin{.}\mathcal Z\arrow[dash,to=1-3,"{\substack{(\mathbb Z/h\mathbb Z)^*\\\unsure{?}}}"]\arrow[dash,to=2-2] \& N(\mu)\arrow[dash,to=1-4]\arrow[dash,to=2-3] \& N\arrow[dash,to=2-4] \&  \\
\mu\arrow[dash,to=2-2]\arrow[dash,to=3-1] \& \mu\times\mathcal Z\arrow[dash,to=2-3,"{(\mathbb Z/h\mathbb Z)^*}"]\arrow[dash,to=3-2] \& N(\mu)\cap N'\arrow[equals,to=2-4] \& N\cap N'\arrow[dash,to=2-5] \& N' \\
0\arrow[dash,to=3-2] \& \mathcal Z \&  \&  \& 
\end{tikzcd}}}

On a une extension can. de $(\mathbb Z/h\mathbb Z)^*$ par $\mu\times\mathcal Z$, les deux extensions \unsure{facteurs|induites} de $(\mathbb Z/h\mathbb Z)^*$ par $\mu$ et $(\mathbb Z/h\mathbb Z)^*$ par $\mathcal Z$ sont contenues \unsure{respectiv$^{\mathrm{mt}}$|respectivement|canoniquement} dans $W'$ et dans $W$. \unsure{Celà|Cela} fait dans un gpe de Weyl des sous-groupes remarquables, définis mod conjugaison.
% archive sheet number: 28
\scanpage{30}
\mspage{\handwritten{8}}
comme ci-dessus la fonctorialité de $\handwritten{\alpha}$ en $P$ et en $Q$, on est donc ramené au cas où on a \deleted{$\unsure{P=\underline{\mathrm Z}/n}$}
\[
P=Q=(\underline{\mathrm Z}/n\underline{\mathrm Z})_{\underline{T}}\ ,
\]
$p$ et $q$ étant la section unité de $(\underline{\mathrm Z}/n\underline{\mathrm Z})_{\underline{T}}$, et on est donc ramené à déterminer l'homomorphisme
\[
\handwritten{\alpha}:(\underline{\mathrm Z}/n\underline{\mathrm Z})_{\underline{T}}\otimes(\underline{\mathrm Z}/n\underline{\mathrm Z})_{\underline{T}}\handwritten{\longrightarrow}\underline{\mathrm{Tor}}_1^{\underline{\mathrm Z}}((\underline{\mathrm Z}/n\underline{\mathrm Z})_{\underline{T}},(\underline{\mathrm Z}/n\underline{\mathrm Z})_T)\ .
\]
La description 2.1.7 nous montre que c'est l'isomorphisme
\[
(\underline{\mathrm Z}/n\underline{\mathrm Z})_{\handwritten{\underline T}}\handwritten{\simeq}\underline{\mathrm{Tor}}_1^{\underline{\mathrm Z}}((\underline{\mathrm Z}/n\underline{\mathrm Z})_{\underline{T}},(\underline{\unsure{\mathbb Z|\mathrm Z}}/n\underline{\mathrm Z})_T)\tag*{(2.1.8.1)}
\]
déduit de la résolution plate canonique
\[
0\handwritten{\longrightarrow}\underline{\mathrm Z}_T\mathrel{\mathop{\handwritten{\longrightarrow}}\limits^{n}}\underline{\mathrm Z}_T\handwritten{\longrightarrow}0\tag*{\handwritten{(2.1.8.2)}}
\]
% circled
\margin{\handwritten{$(\mathbb Z/n\mathbb Z)_{\underline T}$}}
du \uline{premier} facteur,\deleted{dans} \inserted{$P=(\underline{\mathrm Z}/n\underline{\mathrm Z})_{\handwritten{T}},$} qui figure dans le $\underline{\mathrm{Tor}}_1$ .

2.1.9. La définition donnée de (2.1.1) n'est pas symétrique, et on peut, en échangeant les rôles de $P$ et de $Q$ dans la construction précédente, définir un homomorphisme analogue
\[
\handwritten{\beta^n_{P,Q}}:\ {}_nP\otimes{}_nQ\handwritten{\longrightarrow}\underline{\mathrm{Tor}}_1(P,Q)\ ,\tag*{(2.1.9.1)}
\]
en se ramenant d'abord au cas $nP=0$, et utilisant une résolution plate \uline{de} $Q$. On peut aussi définir $\handwritten{\beta}$ en termes de $\handwritten{\alpha}$, de façon équivalente, par la commutativité du diagramme
\diagram{\begin{tikzcd}[ampersand replacement=\&,column sep=large,row sep=large]
\deleted{\text{(2}} \& {}_nP\otimes{}_nQ\handwritten{\arrow[r,"\sim"]}\handwritten{\arrow[d,"{\beta^n_{P,Q}}"']} \& {}_nQ\otimes{}_nP\handwritten{\arrow[d,"{\alpha^n_{Q,P}}"]}\\
\text{(2.1.9.2)} \& \underline{\mathrm{Tor}}_1(P,Q)\handwritten{\arrow[r,"\sim"]} \& \underline{\mathrm{Tor}}_1(\underset{\deleted{\unsure{n}}}{Q},P)\ ,
\end{tikzcd}}
où les flèches horizontales sont les isomorphismes de symétrie. Ceci posé, les relations entre $\handwritten{\alpha}$ et $\handwritten{\beta}$ sont\handwritten{/}données par la formule:
\[
\handwritten{\beta}^{n}_{P,Q}=-{\handwritten{\alpha}}^{n}_{P,Q}\ ,\tag*{\(\substack{\text{(2.1.9.3)}\\\deleted{\text{\unsure{(2.1}}}}\)}
\]
en d'autres termes:
\margin{\handwritten{\begin{tikzpicture}[baseline]\draw (0,1.1) -- (0,0) -- (.11,0);\end{tikzpicture}}}

\uline{Lemme} 2.1.10. Soient $P$ et $Q$ deux \deleted{\unsure{g|C}}Groupes commutatifs, $n\handwritten{>}0$ un entier. Alors le diagramme suivant est \uline{anticommutatif} , où $\handwritten{\alpha}$ est défini dans (2.1.1)\handwritten{:}
% archive sheet number: 29
\scanpage{31}
\margin{Car si $g\in\unsure{\mathcal Z|Z|\mathfrak z}(\gamma)$, $g.D=D'$ est conjuguée de $D$ par élément de $T$, donc mod \unsure{ce|cet} dernier $g$ est dans $\mathcal Z(\gamma)\cap T'$}
\diagram{\resizebox{\linewidth}{!}{\begin{tikzcd}[ampersand replacement=\&,row sep=1.4em,column sep=0.8em,row sep=1.9em,column sep=2.3em]
T\arrow[dash,to=1-2,"{\mathcal Z}"]\arrow[dash,to=2-1] \& \mathcal Z(\gamma)\arrow[dash,to=1-3]\arrow[dash,to=2-2] \& N(T)\arrow[dash,to=1-4]\arrow[dash,to=2-3] \& G\arrow[dash,to=2-4] \\
\gamma\in T\cap N(T')\arrow[dash,to=2-2,"{\mathcal Z}"]\arrow[dash,to=3-1,"{\mathbb Z/h\mathbb Z}"] \& \mathcal Z(\gamma)\cap N(T')={N'}^{\gamma}\arrow[dash,to=2-3]\arrow[dash,to=3-2,"{\mathbb Z/h\mathbb Z}"] \& N(T)\cap N(T')\arrow[dash,to=2-4]\arrow[dash,to=3-3] \& N(T')\arrow[dash,to=3-4,"{W'}"] \\
T\cap T'\arrow[dash,to=3-2,"{\mathcal Z}"]\arrow[equals,to=4-1] \& \mathcal Z(\gamma)\cap T'={T'}^\gamma\arrow[equals,to=3-3]\arrow[phantom,to=4-2,"{\rotatebox{90}{$\simeq$}}"] \& N(T)\cap T'\arrow[dash,to=3-4] \& T' \\
0 \& P'/(1-c)P'(1)\arrow[phantom,to=5-2,"{\rotatebox{90}{$\simeq$}}"] \&  \&  \\
 \& \mathcal Z \&  \& 
\end{tikzcd}}}
\margin{Car \unsure{tt|tout} élément de $W$ qui centralise $c$ est une puissance de $c$}
$\uparrow$ Car tout élément de $T$ qui centralise $D$ est \deleted{\unsure{dans}} l'élément unité
\margin{\uline{\unsure{À|Â} prouver}}

Tout élément de $T$ qui normalise $D$ est dans $\unsure{\kappa|k|K}(\mu_h)$, i.e. $h\psi$ (ou, au choix, les $h\alpha_i$) sont \unsure{comb. lin$^{\mathrm{r}}$|comb. lin.|complétées} à coeff. entiers des éléments $\alpha_i+\psi$.

\unsure{Provient du|Revient donc|Prouvons donc} lemme
\[
\det\begin{pmatrix}
n_1+1 & n_2, & n_3 & \cdots & ,\,n_r\\
\unsure{n_2|n_1|n\,x} & n_2+1 & n_3 & \cdots & n_r\\
\cdots & & & &\\
n_r & n_r & & \cdots & n_r+1
\end{pmatrix}=\sum n_i+1
\]
Un élément de $N(T)\cap N(D)$ \deleted{centralise}\inserted{normalise} $D$, et normalise $T$, donc il normalise \deleted{$\unsure{\mu}$} $\mu_h\simeq\mathbb Z/h\mathbb Z\subset T$ qui est formé des éléments de $T$ qui normalisent $D$. Donc
\[
N(T)\cap N(D)=\deleted{T'\cap\mathrm{Norm}(\mu)}\inserted{N(\mu)\cap N(D)}
\quad\bigl(\underset{\substack{\rotatebox{90}{$\simeq$}\\\mathbb Z/h\mathbb Z=\mu_h}}{\mu}\subset T\text{ engendré par }\gamma\bigr)
\]
% archive sheet number: 30
\scanpage{32}
\mspage{\handwritten{11}}
% archive sheet number: 31
\margin{\handwritten{$\unsure{=|--}$}}
\diagram{\begin{tikzcd}[ampersand replacement=\&,column sep=1.5em]
\& \deleted{\unsure{mx|x}}\,P/nP
\handwritten{\arrow[dl,"{x\handwritten{\mapsto}mx}"']}\handwritten{\arrow[dr,"{\handwritten{x\leftrightarrow mx}}"]} \& \\
{}_nP\handwritten{\arrow[dr,"{\deleted{\unsure{x}}\quad x\handwritten{\mapsto x}}"']} \& \&
{}_{\deleted{\unsure{m}}n\handwritten{'}}P\handwritten{\arrow[dl,"{\deleted{\unsure{x}}\ \handwritten{x\mapsto x}}"]}\\
\& {}_{n\handwritten{'}}P \&
\end{tikzcd}}
\deleted{\unsure{ce}} qui \deleted{montre que le diagramme envisagé} est \inserted{\handwritten{\unsure{evidemment|évidemment}}} commutatif.
\margin{\handwritten{$\big|$}}

\deleted{\unsure{2.2.} Soient maintenant}\uline{Corollaire} 2.1\deleted{\unsure{1|x}}.12. Avec les notations de 2.1.11, on a commutativité dans le diagramme suivant
\diagram{\begin{tikzcd}[ampersand replacement=\&,column sep=large,row sep=large]
\& {}_{n\handwritten{'}}P\otimes{}_{n\handwritten{'}}Q
\handwritten{\arrow[r,"{\substack{\handwritten{m\otimes m}\\\deleted{\unsure{p_m\otimes q_m}}}}"]}
\handwritten{\arrow[d,"{\alpha^{n'}_{P,Q}}"']}
\& {}_nP\otimes{}_nQ\handwritten{\arrow[d,"{\alpha^{n}_{P,Q}}"]}\\
\text{(2.1.12.1)} \& {}_{n\handwritten{'}}\underline{\mathrm{Tor}}_1(P,Q)\handwritten{\arrow[r,"{\substack{\handwritten{m}\\\deleted{\unsure{\underline{\mathrm{Tor}}_1(p,q)|\underline{\mathrm{Tor}}_1(P,Q)}}}}"]} \& {}_n\underline{\mathrm{Tor}}_1(P,Q)\ ,
\end{tikzcd}}
\margin{\handwritten{$\xrightarrow{m}$}}
\deleted{où $q_m:{}_{n'}Q\unsure{\longrightarrow|-}{}_nQ$ est \unsure{dééfini|défini} comme $p_m$ comme induit par $m.\mathrm{id}_Q$ . En}
\margin{\handwritten{$\Bigg|$}}
où, dans la description des flèches horizontales, $m$ désigne les homomorphismes $\mathord{}_{n\handwritten{'}}E\handwritten{\longrightarrow}{}_nE$ induits par $\unsure{m|\bar m|\dot m}.\mathrm{id}_E$ (pour $E$ égal respectivement à $P,Q$ et à $\underline{\mathrm{Tor}}_1(P,Q)$).

En effet, en vertu de \deleted{(}2.1.11 on aura, pour $x\handwritten{\in}{}_{n\handwritten{'}}P(S)$, $y\handwritten{\in}{}_{n\handwritten{'}}Q(S)$: $\handwritten{\alpha}^n(mx\handwritten{\otimes}my)=\handwritten{\alpha}^{n\handwritten{'}}(x\handwritten{\otimes}my)$ , or le deuxième membre est égal à $m\handwritten{\alpha}^{n\handwritten{'}}(x\handwritten{\otimes}y)$ p\deleted{\unsure{p}}uisque $x\otimes my=\deleted{\unsure{mx\otimes}}\ m(x\otimes y)$, d'où la commutativité de (2.1.12.1).

2.1.13. Soit $\unsure{l|\ell}$ un nombre premier. Pour tout Groupe commutatif $E$, désignons par $T_l(E)$ le système projectif \deleted{formé des}
\[
T_l(E)=({}_{l^{\handwritten{\nu}}}E)_{\handwritten{\nu\geq}0}\ ,\tag*{(2.1.13.1)}
\]
qui dépend \unsure{évidemment|évidement} fonctoriellement de $E$. Alors les compatibilités 2.1.12, pour $n,n\handwritten{'}$ de la forme $l^{\handwritten{\nu}},l^{\handwritten{\nu'}}$ avec $\handwritten{\nu'\geq\nu}$, nous montrent que les
\margin{\handwritten{$\alpha^{l^\nu}_{P,Q}$}}
\deleted{$\unsure{\alpha|h}$}\inserted{\handwritten{$\alpha^{l^\nu}_{P,Q}$}} pour $\handwritten{\nu}$ variable définissent un homomorphisme de systèmes projectifs\handwritten{,} \inserted{\handwritten{fonctoriel en $P,Q$:}}
\[
\handwritten{\alpha}^{(l)}_{P,Q}\handwritten{:}\ T_l(P)\otimes T_l(Q)\handwritten{\longrightarrow}T_l(\underline{\mathrm{Tor}}_1(P,Q))\ .\tag*{(2.1.13.2)}
\]
La connaissance de ces homomorphismes, pour tout nombre premier $l$, équivaut évidemment à la connaissance des $\handwritten{\alpha}_{P,Q}^n$ pour tout entier $n$. L\handwritten{es}
\scanpage{34}
Quelques rappels de résultats de Kostant:

a) Pour $\mu_h\subset(T)\subset G$ fixé, les droites cycliques \inserted{$D$} $1$-normalisées par $\mu_h$ (\inserted{i.e.,} avec car $1$) forment un espace homogène sous $T$, stabilisateur $\mu_h$, donc un espace homogène sous $\mu_h$.
Les droites cycliques $D$ correspondent biunivoquement aux sous tores $T'$ de $G$ \deleted{tels que}\inserted{en opposition} \deleted{$\unsure{T|\mu_h}$} à $T$ pour $\mu_h$ (i.e. normalisé par $\mu_h$, et tel que $\mu_h$ y admet un \deleted{\unsure{elt}}\inserted{car.} propre au moins qui est \deleted{\unsure{racine}} car. \uline{primitif}, ($\unsure{\simeq|\Leftrightarrow}$ le car. $1$ y est car propre): à $D$ correspond $T'=\mathrm{Cent}(D)$, à $T'$ correspond l'\deleted{unique droite de}\unsure{le sous-espace|le sous-espace propre|unique sous-espace} propre de $\mu_h$ opérant sur \deleted{$\unsure{\mathfrak t}$} $\mathrm{Lie}\,T'\subset\mathfrak g$. On a d'ailleurs $T\cap N(T')=\mu_h$.

b) Les droites cycliques $D$ de $\mathfrak g$ \uline{mod} $\mu_h^*$ s'identifient can. \unsure{avec|à|via} une droite $\Delta$ canonique sur $k$. On l'obtient par descente à partir de $D^{\otimes h}$ défini sur le schéma des droites cycliques qui devient un schéma en droites sur le schéma des tores maximaux $T'$, et à ce titre se descend.
% slanted
\margin{NB Si on se place sur un topos, tel le classifiant ${\unsure{\underline B|B}}_{E_k}^{\deleted{\unsure{*}}}$, on ne pourra pas en général trouver de $a$, donc on ne pourra définir de torseur $\mathrm{KostKillnorm}$: cela provient du fait que la repr. de $E$ sur $\Delta$ peut être non triviale. Il faut passer à \unsure{rev.|une} étale d'indice $2$ \unsure{---}. De façon générale, on a obstruction \deleted{\unsure{pour}}\inserted{à trouver $D$} dans $H^2(\unsure{Q|\mathbb Q};E;\mu_h)$ nulle si on restreint \unsure{à|:} $H^{\unsure{2|1}}(Q,\mu_h)$ ou $H^2(Q;E^\circ;\mu_h)$}
On l'obtient aussi comme l'image du lieu cyclique de $\mathfrak g$ par l'\unsure{hom|hom.} de Kostant $\mathfrak g\longrightarrow I=\mathfrak t/W$, qui a une structure \inserted{\uline{vectorielle}} \deleted{\unsure{linéaire}} car le stabilisateur \inserted{dans $W$} d'une droite cyclique $D\subset\mathfrak t$ y opère par $\mu_h$, et $D/\mu_h$ a une structure \deleted{de droite}\inserted{vectorielle} \ldots

Il serait intéressant de calculer \deleted{les relations de \unsure{cette}} la repr qu'on obtient ainsi de $\unsure{E|\mathcal E|G}$ (autom. ext.) dans $\mathbb Z/2\mathbb Z\simeq\mu_2(\mathbb Z)$, et la comparer avec le car. provenant de la repr de $E$ sur $\mathfrak g_\psi$ ($\psi$ plus grande racine). P. ex. il est clair par les résultats de Kostant que si \deleted{la} deuxième repr. est triviale, la première aussi.
A-t-on, pour tt tore $T$ dans un Borel \deleted{\unsure{g}}$B$, un isom can $\mathfrak g_\psi\simeq\Delta$ (ou $\mathfrak g_\psi\simeq\Delta^{-1}$, \deleted{\unsure{c}} cela revient au $=$, car \unsure{on a} un isom can $\Delta^{\otimes2}\simeq 1$).
% archive sheet number: 33
\scanpage{35}
Considérons le schéma des \deleted{triplets} systèmes $\mathrm{RigKos}$ $(T,T',B,B')$ avec
\[
\begin{aligned}
&T,T' \text{ tores maximaux en opposition}\\
&B,B' \text{ des Borels contenant }T,T'\text{ respectivement.}
\end{aligned}
\qquad\big[(h,p)=1\big]
\]
\inserted{(\uline{Rigidifications de Kostant})}
C'est un schéma lisse sur $S$, sur lequel $G$ opère librement.
\[
\begin{aligned}
&\mathrm{RigKos}/G \simeq \text{schéma fini étale sur }S,\ \text{\deleted{d'ordre} de rang }{\unsure{w|W}}^2/\varphi(h)hz\\
&\hspace{2em}\left\{\begin{aligned}
&\text{où }w\text{ est l'ordre du gpe de Weyl}\\
&\text{\deleted{\unsure{et}} où }h\text{ est le \unsure{n$^{\mathrm b}$|n°} de Coxeter}\\
&\varphi\text{ l'indicatrice d'Euler}\\
&z\text{ ordre du centre}
\end{aligned}\right.
\end{aligned}
\]
Il est associé à une représentation du gpe $E$ des autom extérieurs dans un schéma \inserted{fini} étale sur $\operatorname{Spec}\mathbb Z[1/h]$ [le schéma des \uline{classes} de rigidifications de Kostant de $G$ sur $\operatorname{Spec}(\mathbb Z[1/h])$].
Le dernier schéma est \deleted{bel et bien}\inserted{il} \deleted{de} splitté : prenons un couple en \unsure{opposit|opposition},  défini par
\[
\big(\underbrace{\mu_h\mathrel{\overset{\kappa}{\text{---}}}T}_{\substack{\text{canonique}\\\text{défini par épinglage}}},\quad
D\text{ engendré par }\ e_{\alpha_1}\unsure{\,\text{---}\,|+\cdots+}e_{\alpha_r}+
\underbrace{e_{-\psi}}_{\substack{\text{défini au}\\\text{signe près}}}=e
\]
On trouve $T'$ comme centralisateur \deleted{\unsure{commun}} de cet élément.
On a \deleted{le gpe} Alors le schéma des Borels \inserted{$\unsure{Q|\text{il}}$} contenant $T$, isomorphe à $W_{\mathbb Z}$, et le schéma $Q'$ des Borels contenant $T'$, qui est très tordu, \deleted{\illegible{1}} c'est un torseur sous $W'=N(T')/T'$.
On a dans $W\times W'$ le sous groupe \inserted{$\Phi=N\cap N'$} extension de $(\mathbb Z/h\mathbb Z)^*$ par $\deleted{\unsure{\mu|\mu_h}}\,\unsure{\mathcal Z|Z|\mathfrak z}\times\mu_h$, où $\mathcal Z\subset W$ (centralisateur de $\kappa$) et $\mu_h\subset W'$. Il opère \inserted{librement} sur $Q\times Q'$ à travers $W\times W'$, et on forme $(Q\times Q')/\Phi$.
Le gpe des autom extérieurs \inserted{$E$} de $G$ y opère via son action par $\pm1$ sur $\mathfrak g_{-\psi}$, donc en fait, \deleted{\unsure{ou celle}} ou bien l'opération de $E$ est triviale, ou elle se fait par une autom can.
% archive sheet number: 34
\scanpage{36}
\mspage{\handwritten{6}}
2.\uline{Accouplements définis par une biextension}.
2.1. Soient $P$ et $Q$ deux \deleted{g}Groupes abéliens sur $T$, et $n\handwritten{>}0$ un entier. Rappelons qu'on désigne par $\mathord{}_nP$ l\deleted{a}e noyau de la multiplication par $n\,\mathrm{id}_P$. Nous allons définir un homomorphisme canonique
\[ \handwritten{\alpha}=\handwritten{\alpha}^{\deleted{\unsure{N|n}}n}_{P,Q}:\quad {}_nP\otimes{}_nQ\mathbin{\handwritten{\longrightarrow}}\underline{\mathrm{Tor}}_1(P,Q),\tag*{(2.1.1)} \]
qui sera le composé
\[ {}_nP\otimes{}_nQ\mathbin{\handwritten{\longrightarrow}}\underline{\mathrm{Tor}}_1({}_nP,{}_nQ)\mathbin{\handwritten{\longrightarrow}}\underline{\mathrm{Tor}}_1(P,Q),\tag*{(2.1.2)} \]
où la première flèche est $\handwritten{\alpha}^n_{{}_nP,{}_nQ}$, et la deuxième est celle déduite des inclusions $\mathord{}_nP\handwritten{\to}P$, $\mathord{}_nQ\handwritten{\to}Q$. On est donc ramené à définir (2.1.1) dans le cas où on a $nP=0,nQ=0$, i.e. où $P$ et $Q$ sont des $\underline{\mathrm Z}/n\underline{\mathrm Z}$ - Modules.
\deleted{Calculant alors $\underline{\mathrm{Tor}}_{\unsure{1|i}}(P,Q)$ à l'aide d'une} Supposons seulement \deleted{\unsure{$n$|$nP$}}\inserted{$nQ=0$.} Prenant une résolution plate $\deleted{\unsure{M}}L\mathord.$ de $\deleted{Q}P$
\deleted{\illegible{1}}
(en tant que $\underline{\mathrm Z}$-Module), \deleted{comme $\underline{\mathrm H}_{\unsure{1}}(P\otimes M\mathord.)$ \unsure{$=$}} on trouve des isomorphismes canoniques
\[ \deleted{\mathrm H\mathord.(L\mathord.\otimes Q)=\mathrm H\mathord.((L\mathord.\otimes_{\underline{\mathrm Z}}}\qquad\deleted{P\otimes} \]
\[ P\overset{\underline L}{\otimes}_{\underline{\mathrm Z}}Q\handwritten{\simeq}L\mathord.\otimes_{\underline{\mathrm Z}}Q\handwritten{\simeq}(L\mathord.\otimes_{\underline{\mathrm Z}}\underline{\mathrm Z}/n\underline{\mathrm Z})\otimes_{\underline{\mathrm Z}/n\underline{\mathrm Z}}Q, \]
qu'on peut écrire comme un isomorphisme dans la catégorie dérivée
\[ P\overset{\underline L}{\otimes}_{\underline{\mathrm Z}}Q\handwritten{\simeq}(P\overset{\underline L}{\otimes}_{\underline{\mathrm Z}}\underline{\mathrm Z}/n\underline{\mathrm Z})\overset{\underline L}{\otimes}_{\underline{\mathrm Z}/n\underline{\mathrm Z}}Q,\tag*{(2.1.3)} \]
donnant naissance à \deleted{\unsure{une}} \handwritten{la} suite spectrale de type homologique \inserted{\handwritten{“de Künneth”}} : \margin{\handwritten{(cf. SGA 1 $\overline{\underline{\mathrm{IV}}}$ 5.2)}}
\[ \underline{\mathrm{Tor}}^{\underline{\mathrm Z}}_*(P,Q)\handwritten{\Longleftarrow}E^2_{p,q}=\underline{\mathrm{Tor}}^{\underline{\mathrm Z}/n\underline{\mathrm Z}}_{\deleted{\unsure{q}}p}(\underline{\mathrm{Tor}}^{\underline{\mathrm Z}}_q(P,\underline{\mathrm Z}/n\underline{\mathrm Z}),Q).\tag*{(2.1.4)} \]
Rappelons d'ailleurs qu'on a des isomorphismes canoniques
\[ \underline{\mathrm{Tor}}^{\underline{\mathrm Z}}_0(P,\underline{\mathrm Z}/n\underline{\mathrm Z})=P_n\mathrel{\overset{\mathrm{dfn}}{=}}P/nP,\quad \underline{\mathrm{Tor}}^{\underline{\mathrm Z}}_1(P,\underline{\mathrm Z}/n\underline{\mathrm Z})\simeq {}_nP,\quad \underline{\mathrm{Tor}}^{\underline{\mathrm Z}}_i(P,\underline{\mathrm Z}/n\underline{\mathrm Z})=0\text{ si }i\geq2, \]
déduits du calcul des $\underline{\mathrm{Tor}}^{\underline{\mathrm Z}}_i(P,\underline{\mathrm Z}/n\underline{\mathrm Z})$ à l'aide de la résolution plate canonique suivante de $(\underline{\mathrm Z}/n\underline{\mathrm Z})_T$ :
\[ \begin{aligned}0\deleted{\unsure{n}}&\longrightarrow\underline{\mathrm Z}_T\overset{n}{\longrightarrow}\underline{\mathrm Z}_T\longrightarrow 0\deleted{\unsure{1}}\\ &\hspace{2.6em}\overset{\prime\prime}{\underline L_1}\qquad\overset{\prime\prime}{\underline L_0}\qquad,\unsure{\prime\prime}\end{aligned} \]
l'augmentation $\underline L_0\to(\underline{\mathrm Z}/n\underline{\mathrm Z})_T$ étant définie en envoyant 1 sur 1. Par suite, la suite spectrale (2.1.4) se réduit essentiellement à une suite exacte
\[ 0\to\underline{\mathrm{Tor}}^{\underline{\mathrm Z}/n\underline{\mathrm Z}}_2(P_n,\deleted{\unsure{*}}Q)\to{}_nP\otimes Q\mathbin{\overset{\handwritten{\alpha}}{\handwritten{\longrightarrow}}}\underline{\mathrm{Tor}}^{\underline{\mathrm Z}}_1(P,Q)\to\underline{\mathrm{Tor}}^{\underline{\mathrm Z}/n\underline{\mathrm Z}}_1(P_n,Q)\to0,\tag*{(2.1.5)} \]
où on a posé $P_n=P\otimes_{\underline{\mathrm Z}}\underline{\mathrm Z}/n\underline{\mathrm Z}=P/nP$. L'homomorphisme $\handwritten{\alpha}$ \deleted{\unsure{est}} cherché est \deleted{celui} \inserted{\handwritten{l'homomorphisme-coin $E^2_{0,1}\to\underline{\mathrm{Tor}}^{\mathbb Z}_1$}} qui figure dans la dernière suite exacte. Cette dernière étant évidemment fonctorielle en $P$ et en $Q$, la fonctorialité en $P$ montre que l'homomorphisme envisagé est bien le composé
\[ \handwritten{{}_n}P\otimes Q\to\underline{\mathrm{Tor}}_1\handwritten{({}_nP,Q)\longrightarrow\underline{\mathrm{Tor}}_1(P,Q),}\qquad\deleted{\unsure{N|n}} \]
% archive sheet number: 35
